% GENERATED by tools/sizing-model/camera_optics_tex.py -- do not edit.
\documentclass[10pt,a4paper]{article}
\usepackage[margin=18mm]{geometry}
\usepackage{amsmath,booktabs,array,xcolor}
\definecolor{hi}{HTML}{0B5394}
\definecolor{dim}{HTML}{5A5A62}
\definecolor{warn}{HTML}{9C0006}
\setlength{\parindent}{0pt}
\setlength{\parskip}{4pt}
\pagestyle{empty}
\renewcommand{\arraystretch}{1.25}

\begin{document}

{\LARGE\bfseries Camera and lens: complete derivation}\\[1mm]
{\color{dim}Project RescueSwarm $\cdot$ NIDAR 2026--27 Track 1 $\cdot$
generated from \texttt{tools/sizing-model/camera\_optics.py}}\\[1mm]
{\color{hi}\rule{\linewidth}{1pt}}

\section*{1. The part, and a discrepancy}

The bill of materials buys an \textbf{Arducam IMX477} with a 6\,mm S-mount
lens. Sony specifies the IMX477 as type 1/2.3: $4056\times3040$ at a
\textbf{1.55\,\textmu m} pixel pitch, 7.9\,mm diagonal.

\textcolor{warn}{\textbf{The sizing chapter was not written against that
sensor.}} \texttt{sizing-calculations.md}~\S8 states the baseline as
``1/1.8\,in, $4056\times3040$, $7.4\times5.6$\,mm, \textbf{1.82\,\textmu m}
pitch''. The pixel \emph{count} matches and the pixel \emph{size} does not.
Every optical quantity below scales with pitch, so none of the published
figures describe the camera being bought. Both cases are computed here.

\section*{2. Sensor geometry}

Active area follows from count and pitch:
\[
w = N_x\,p, \qquad h = N_y\,p, \qquad d = \sqrt{w^2+h^2}
\]
\[
w = 4056 \times 1.55\,\text{\textmu m} = 6.287\,\text{mm},
\quad
h = 3040 \times 1.55\,\text{\textmu m} = 4.712\,\text{mm},
\quad
d = 7.86\,\text{mm}
\]

\section*{3. Field of view}

The lens forms the image at its focal length, so each half-field is the
arctangent of half the sensor over $f$ --- exact, with no small-angle
approximation:
\[
\theta = 2\arctan\!\left(\frac{s}{2f}\right)
\]
\[
\text{HFOV} = 2\arctan\!\frac{6.287}{12.0}
= 55.3^\circ,
\qquad
\text{VFOV} = 42.9^\circ,
\qquad
\text{DFOV} = 66.4^\circ
\]
One pixel subtends 14.80\,mdeg, which is the quantity that sets
the blur tolerance in~\S6.

\section*{4. Ground sample distance and footprint}

Sensor and ground form similar triangles about the lens, so
\[
\text{GSD} = \frac{p\,H}{f},
\qquad
W = 2H\tan\!\frac{\text{HFOV}}{2},
\qquad
S = W\,(1-\ell)
\]
for altitude $H$, pitch $p$, focal length $f$ and sidelap
$\ell = 30\,\%$.
Sidelap exists because attitude wobble, altitude error and lens distortion all
move the real footprint around; at zero sidelap a gap between passes is an
undetected survivor.

\begin{center}
\begin{tabular}{r r r c r r r}
\toprule
\textbf{AGL} & \textbf{GSD} & \textbf{Survivor} & \textbf{Footprint} &
\textbf{Spacing} & \textbf{Lines} & \textbf{Sweep} \\
(m) & (cm/px) & (px) & (m) & (m) & & (s) \\
\midrule
30 & 0.78 & 219 & 31.4 $\times$ 23.6 & 22.0 & 9 & 253 \\
40 & 1.03 & 165 & 41.9 $\times$ 31.4 & 29.3 & 7 & 196 \\
60 & 1.55 & 110 & 62.9 $\times$ 47.1 & 44.0 & 5 & 138 \\
\bottomrule
\end{tabular}
\end{center}

Transect count and sweep time assume the 10\,ha area split
3 ways, each sub-region taken as square, flown at
8\,m/s with a 6\,s cost per turn.

\section*{5. What the substitution changes, at 40\,m}

\begin{center}
\begin{tabular}{l r r r}
\toprule
 & \textbf{IMX477} & \textbf{Sizing model} & \textbf{Change} \\
\midrule
Horizontal field of view & 55.30 & 63.20 & -12.5\,\% \\
Ground sample distance & 1.03 & 1.21 & -14.8\,\% \\
Survivor, long axis & 164.52 & 140.11 & +17.4\,\% \\
Swath width & 41.91 & 49.21 & -14.8\,\% \\
Line spacing & 29.34 & 34.45 & -14.8\,\% \\
Transects per drone & 7.00 & 6.00 & +16.7\,\% \\
Sweep time per drone & 195.75 & 166.93 & +17.3\,\% \\
\bottomrule
\end{tabular}
\end{center}

A smaller pixel behind the same lens narrows the field and finens the ground
sample distance in the same proportion. \textbf{Detection improves by
17\,\% and coverage costs
17\,\% more sweep time.} Against 250
marks for detection and a 1800\,s mission budget, that trade favours the part
actually specified --- but it was never chosen, and the published figures
describe neither camera.

\section*{6. Motion blur}

Forward motion smears the image by $v\,t_{\text{exp}}$ on the ground. Holding
that under one pixel:
\[
t_{\text{exp}} \le \frac{\text{GSD}}{v}
\]

\begin{center}
\begin{tabular}{r r r}
\toprule
\textbf{AGL} (m) & \textbf{Max exposure} (ms) & \textbf{Shutter} \\
\midrule
30 & 0.97 & 1/1032 \\
40 & 1.29 & 1/774 \\
60 & 1.94 & 1/516 \\
\bottomrule
\end{tabular}
\end{center}

\texttt{SYS-45} already requires $\le 1/1000$\,s, which clears every altitude
with margin. Blur is not a binding constraint; it is Indian daylight driving
ISO down rather than exposure up.

\section*{7. Tiling}

A survivor occupies a few tens of pixels in a 12\,MP frame, so the detector is
fed overlapping crops rather than the whole image. Tile count depends on pixel
\emph{count}, so it is unaffected by the sensor substitution --- but the
\emph{effective} GSD after downsampling is not.

\begin{center}
\begin{tabular}{c c r r r r}
\toprule
\textbf{Downsample} & \textbf{Frame} & \textbf{Tiles} &
\textbf{Inf./s} & \textbf{GSD} (cm/px) & \textbf{Survivor} (px) \\
\midrule
$1\times$ & 4056$\times$3040 & 48 & 240 & 1.03 & 165 \\
$2\times$ & 2028$\times$1520 & 12 & 24 & 2.07 & 82 \\
\bottomrule
\end{tabular}
\end{center}

At $2\times$ downsampling and 2\,Hz the budget closes at 24 inferences per
second, unchanged. The survivor is
\textbf{82\,px} in that downsampled frame
--- not the 47\,px the proposal currently states, which was computed at 60\,m on
the other sensor.

\section*{8. What follows}

Nothing here argues for changing the camera. The IMX477 is the better of the
two for detection, which is where the marks are. What it argues is that the
\emph{documents} should be re-baselined onto the part being bought: the GSD
table, the swath and line-spacing figures, the sweep times and the survivor
pixel count in the perception section. That work interacts with the unresolved
40\,m versus 60\,m survey-altitude question, because both move the same
quantities, and the two should be settled in one pass rather than two.

\vfill
{\footnotesize\color{dim}Generated file. Regenerate with\\
\texttt{python tools/sizing-model/camera\_optics\_tex.py} then
\texttt{pdflatex}.\\
Every figure above is computed by
\texttt{tools/sizing-model/camera\_optics.py}.}

\end{document}
